\begin{theorem}[Exact elliptic trace formula for the toric fibers]
\label{thm:toric-fiber-sym2}
Let $p\geq5$ be prime and $a\in\mathbb F_p^\times$.  Put
\[
 D_a=a^2-34a+1,
 \qquad \epsilon_a=\mathbf 1_{a=1},
\]
and let $\theta_p(a)$ be the primitive weight-two trace in
Theorem~\ref{thm:toric-fiber-K3}.  On the quadratic cover
\begin{equation}
 8x^2+(a-1)x+a=0
 \label{eq:toric-sym2-cover}
\end{equation}
consider the elliptic curve
\begin{equation}
 E_x:\quad Y^2+(1-2x)XY+x^2Y=X^3.
 \label{eq:toric-sym2-elliptic-curve}
\end{equation}
For an elliptic curve $E/\mathbb F_q$, write
\[
 A_q(E)=q+1-\#E(\mathbb F_q),
\]
and write $\chi_p$ for the quadratic character of~$\mathbb F_p$.
Then the following integer identities hold.

\begin{enumerate}
\item If $D_a$ is a nonzero square and $x\in\mathbb F_p$ is either
root of~\eqref{eq:toric-sym2-cover}, then
\begin{equation}
 \boxed{\theta_p(a)=A_p(E_x)^2-p-\epsilon_a p.}
 \label{eq:toric-sym2-split}
\end{equation}

\item If $D_a$ is a nonsquare and $x\in\mathbb F_{p^2}$ is either
root of~\eqref{eq:toric-sym2-cover}, then
\begin{equation}
 \boxed{\theta_p(a)=
 \chi_p(-3)A_{p^2}(E_x)-p-\epsilon_a p.}
 \label{eq:toric-sym2-inert}
\end{equation}

\item If $D_a=0$ and $x\in\mathbb F_p$ is the double root, then
\begin{equation}
 \boxed{\theta_p(a)=
 A_p(E_x)^2-p-\chi_p(-6)p.}
 \label{eq:toric-sym2-branch}
\end{equation}
\end{enumerate}
In particular, each right-hand side is independent of the chosen root, and
\begin{equation}
 \mu_p(a)=p^2-6p+12+\theta_p(a)
 \label{eq:toric-sym2-mu}
\end{equation}
is an exact elliptic-curve expression for every nonzero toric fiber.
\end{theorem}

\begin{proof}
We first identify the surface and its Shioda--Inose cover algebraically.  The
Beukers--Peters affine model~\cite{BeukersPeters1984} is
\begin{equation}
 1-(1-\mathsf X\mathsf Y)\mathsf Z
 -a\mathsf X\mathsf Y\mathsf Z
       (1-\mathsf X)(1-\mathsf Y)(1-\mathsf Z)=0.
 \label{eq:toric-sym2-BP-model}
\end{equation}
Under the birational change of variables
\begin{equation}
 \mathsf X=\frac{y}{1+y},\qquad
 \mathsf Y=\frac{z}{1+z},\qquad
 \mathsf Z=\frac{x}{1+x},
 \label{eq:toric-sym2-BP-change}
\end{equation}
one has the rational-function identity
\begin{equation}
 \begin{aligned}
 \frac{1-(1-\mathsf X\mathsf Y)\mathsf Z}
 {\mathsf X\mathsf Y\mathsf Z
  (1-\mathsf X)(1-\mathsf Y)(1-\mathsf Z)}
 &=\frac{(1+x)(1+y)(1+z)}{xyz}\\
 &\hspace{2em}\cdot\bigl((1+y)(1+z)+xyz\bigr)
 =\Lambda(x,y,z).
 \end{aligned}
 \label{eq:toric-sym2-BP-identity}
\end{equation}
Thus~\eqref{eq:toric-sym2-BP-model} is precisely the fiber
$\Lambda=a$, and hence is birational to the surface in
\eqref{eq:toric-K3-affine}; this is stronger than a comparison of
Picard--Fuchs equations.

Peters--Stienstra~\cite{PetersStienstra1989} identify the base change
$a=s^2$ of this pencil with the Ap\'ery--Fermi pencil.  The
Morrison--Nikulin quotient of fibration~$\#16$ in
Bertin--Lecacheux~\cite[Corollary~4.1]{BertinLecacheux2020} is the Kummer
surface of a pair of $6$-isogenous elliptic curves.  Their displayed rational
maps are defined over the corresponding quadratic cover, so their graphs give
the Shioda--Inose correspondence in $\ell$-adic cohomology after reduction at
every $p\geq5$.  We now identify the two elliptic factors, including the
arithmetic twist.

The deck involution of~\eqref{eq:toric-sym2-cover} is
\begin{equation}
 \sigma(x)=\frac{1-8x}{8(1+x)},
 \qquad
 \frac{x(1-8x)}{1+x}
 =\frac{\sigma(x)(1-8\sigma(x))}{1+\sigma(x)}=a.
 \label{eq:toric-sym2-deck}
\end{equation}
The two standard symmetric functions of the modular invariants are
\begin{align}
 j(E_x)j(E_{\sigma(x)})
 &=\frac{(a+1)^3}{a^7}
   (a^3-21a^2+219a+1)^3,
 \label{eq:toric-sym2-j-product}\\
 (j(E_x)-1728)(j(E_{\sigma(x)})-1728)
 &=\frac{(a^2+14a+1)^2}{a^7}\notag\\
 &\hspace{2em}\cdot
 (a^4-44a^3+198a^2-548a+1)^2.
 \label{eq:toric-sym2-j-shifted-product}
\end{align}
These are exactly the two formulas for the elliptic factors of the
Morrison--Nikulin quotient, after the harmless choice $s^2=a^{-1}$;
the Fermi parameter $s+s^{-1}$ is unchanged by $s\mapsto s^{-1}$.
Equations~\eqref{eq:toric-sym2-j-product} and
\eqref{eq:toric-sym2-j-shifted-product} determine the unordered pair of
$j$-invariants, since they determine both its product and its sum.

The twist can be fixed without an appeal to periods.  On $E_x$ the point
\[
 Q=(x,-x)
\]
has exact order six; its successive nonzero multiples have coordinates
\begin{align*}
 Q&=(x,-x),& 2Q&=(0,-x^2),&3Q&=(-x^2,-x^3),\\
 4Q&=(0,0),&5Q&=(x,x^2).&&
\end{align*}
V\'elu's formulas give the quotient $C_x=E_x/\langle Q\rangle$ as
\begin{align}
 C_x:\quad Y^2+(1-2x)XY+x^2Y
 ={}&X^3+c_4^{\flat}(x)X+c_6^{\flat}(x),
 \label{eq:toric-sym2-Velu-quotient}\\
 c_4^{\flat}(x)={}&-5x^4-5x^3-20x^2-5x,\notag\\
 c_6^{\flat}(x)={}&3x^6+7x^5-41x^4-24x^3-14x^2-x.\notag
\end{align}
Here the superscript $\flat$ is only a label; these are the Weierstrass
coefficients $a_4,a_6$, not the usual invariants.  If $\omega^2=-3$, put
\begin{align*}
 q&=\frac{4(1+x)\omega}{3},
 &r&=-\frac{13x^2+5x+1}{3},\\
 s_0&=\frac{4x\omega+2x+\omega-1}{2},
 &t_0&=-\frac{64x^3\omega+78x^3+48x^2\omega
 -15x\omega-9x+\omega-3}{18}.
\end{align*}
The standard Weierstrass change
\[
 X_C=q^2X_{\sigma}+r,qquad
 Y_C=q^3Y_{\sigma}+q^2s_0X_{\sigma}+t_0
\]
identifies $E_{\sigma(x)}$ with $C_x$ over
$\mathbb Q(x,\sqrt{-3})$.  Equivalently,
\begin{align}
 \frac{c_4(C_x)}{c_4(E_{\sigma(x)})}
 &=\frac{256}{9}(1+x)^4,\notag\\
 \frac{c_6(C_x)}{c_6(E_{\sigma(x)})}
 &=-\frac{4096}{27}(1+x)^6,\notag\\
 \frac{\Delta(C_x)}{\Delta(E_{\sigma(x)})}
 &=\frac{2^{24}}{3^6}(1+x)^{12}.
 \label{eq:toric-sym2-twist-invariants}
\end{align}
Thus the twist square class is the constant $-3$.  In particular, whenever
$x,\sigma(x)\in\mathbb F_p$,
\begin{equation}
 A_p(E_{\sigma(x)})=\chi_p(-3)A_p(E_x).
 \label{eq:toric-sym2-trace-twist}
\end{equation}

We recall the resulting cohomological calculation.  Let $V_x$ be
$H^1_{\mathrm{et}}(E_{x,\overline{\mathbb F}_p},\mathbb Q_\ell)$.
For $D_a\ne0$, the rank-four Kummer summand is
\begin{equation}
 \chi_p(-3)\otimes
 \operatorname {As}^{+}_{\mathbb F_p[x]/\mathbb F_p}(V_x).
 \label{eq:toric-sym2-Asai}
\end{equation}
The graph of the six-isogeny contributes a rational Tate line
$\mathbb Q_\ell(-1)$ to this summand.  Its orthogonal complement is the
generic rank-three transcendental part of the K3 surface.  If the quadratic
algebra splits, the trace of~\eqref{eq:toric-sym2-Asai} is
\[
 \chi_p(-3)A_p(E_x)A_p(E_{\sigma(x)})=A_p(E_x)^2
\]
by~\eqref{eq:toric-sym2-trace-twist}.  If it is the field
$\mathbb F_{p^2}$, the elementary tensor-induction trace identity gives
\[
 \operatorname {Tr}(\operatorname {Frob}_p\mid
 \operatorname {As}^{+}V_x)=A_{p^2}(E_x),
\]
because Frobenius interchanges the two tensor factors.  Removing the Tate
line therefore gives respectively
\[
 A_p(E_x)^2-p,
 \qquad
 \chi_p(-3)A_{p^2}(E_x)-p.
\]
At $a=1$ the explicit divisor calculation in
Theorem~\ref{thm:toric-fiber-K3} supplies one additional rational Tate line;
this accounts for the term $-\epsilon_ap$ in
\eqref{eq:toric-sym2-split} and~\eqref{eq:toric-sym2-inert}.

It remains to check the ramification points of the quadratic cover.  There
$\sigma(x)=x$, so
\begin{equation}
 8x^2+16x-1=0,
 \qquad (1+x)^2=\frac98.
 \label{eq:toric-sym2-branch-fixed}
\end{equation}
The six-isogeny followed by the displayed twist isomorphism is now an
endomorphism $\alpha$ of $E_x$.  Since the V\'elu isogeny is normalized, its
action on the invariant differential is multiplication by~$q$, and
\[
 q^2=-\frac{16}{3}(1+x)^2=-6.
\]
Consequently $\alpha^2=[-6]$ and $E_x$ has complex multiplication by the
order of discriminant $-24$.  Equivalently, its modular invariant satisfies
\[
 j(E_x)^2-4834944j(E_x)+14670139392=0.
\]
The trace-zero CM endomorphism contributes an algebraic line, namely
\[
 \mathbb Q_\ell(-1)\otimes\chi_p(-6).
\]
Its Frobenius trace is $\chi_p(-6)p$.  Removing it from the specialized
rank-three trace
$A_p(E_x)^2-p$ proves~\eqref{eq:toric-sym2-branch}.
Finally, substituting the three formulas into
\eqref{eq:toric-K3-deflated-trace} proves
\eqref{eq:toric-sym2-mu}.
\end{proof}

\begin{remark}[What the elliptic formula does not prove]
\label{rem:toric-sym2-Mellin-limit}
Theorem~\ref{thm:toric-fiber-sym2} replaces every primitive K3 trace by an
exact symmetric-square or inert Asai trace.  It does not by itself give
Mellin anti-concentration in the parameter~$a$.  In particular, a bound for
the number of vanishing Ap\'ery coefficients still requires cancellation
among the multiplicative Mellin transforms of these traces; the pointwise
Hasse bound alone cannot provide such cancellation.
\end{remark}
